% ============================================================
%  Hp3 manuscript -- Cell Reports (Cell Press) submission VARIANT
% ------------------------------------------------------------
%  This is a Cell-Reports-flavoured build that REUSES body.tex verbatim.
%  Differences vs manuscript.tex (the Microbial Biotechnology build):
%    1. References  : numbered (Cell style) via natbib[numbers] + unsrtnat
%    2. Front matter : Highlights + In Brief (eTOC blurb) + Summary,
%                      replacing the plain abstract (Cell Reports layout).
%    3. Significance : Summary/Highlights foreground novelty + in-vivo outcome.
%
%  Cell Press allows FLEXIBLE / format-free INITIAL submission, so:
%    - Methods stay as the narrative "Experimental Procedures" in body.tex
%      (kept before Results for now); convert to STAR Methods at revision.
%    - A PDF is accepted at initial submission.
%
%  Build:  latexmk manuscript-cr      (uses .latexmkrc -> xelatex + bibtex)
%  Body :  body.tex   (SHARED with the MBT build -- edit once, both rebuild)
%  Figs :  ../analyses/data/121-figs/figN.png
%  Refs :  bibliography.bib  (symlink -> Zotero export)
%
%  REFS = author-year (Cell Press style) -- reused AS-IS for iScience and any
%  other Cell Press Multi-Journal Submission target (all author-year + STAR
%  Methods), so co-submitting CR + iScience needs only ONE format.
%  FALLBACK (only if the whole Cell Press route fails) -- Communications Biology
%  (Nature Portfolio) uses NUMBERED refs: switch to
%  \usepackage[numbers,sort&compress]{natbib}\bibliographystyle{unsrtnat},
%  delete the Highlights/In Brief blocks, move Methods after Discussion.
% ============================================================
\documentclass[11pt,a4paper]{article}

% ---- engine / fonts (xelatex; Times New Roman) ----
\usepackage{iftex}
\ifPDFTeX
  \usepackage[T1]{fontenc}\usepackage[utf8]{inputenc}\usepackage{textcomp}
\else
  \usepackage{fontspec}
  \defaultfontfeatures{Scale=MatchLowercase}
  \setmainfont{Times New Roman}
\fi

% ---- layout ----
\usepackage[a4paper,margin=1in]{geometry}
\IfFileExists{microtype.sty}{\usepackage{microtype}}{}
\usepackage{setspace}\onehalfspacing
\setlength{\parindent}{0pt}\setlength{\parskip}{6pt plus 2pt minus 1pt}
\setlength{\emergencystretch}{3em}
\setcounter{secnumdepth}{5}

% ---- math / graphics / tables (dependencies of body.tex) ----
\usepackage{amsmath,amssymb}
\usepackage{graphicx}
\usepackage{longtable,booktabs,array,calc,etoolbox}
\usepackage{caption,subcaption}
\usepackage{xcolor}

% ---- pandoc helper macros required by body.tex ----
\makeatletter
\newsavebox\pandoc@box
\newcommand*\pandocbounded[1]{% scale image to fit text height/width
  \sbox\pandoc@box{#1}%
  \Gscale@div\@tempa{\textheight}{\dimexpr\ht\pandoc@box+\dp\pandoc@box\relax}%
  \Gscale@div\@tempb{\linewidth}{\wd\pandoc@box}%
  \ifdim\@tempb\p@<\@tempa\p@\let\@tempa\@tempb\fi
  \ifdim\@tempa\p@<\p@\scalebox{\@tempa}{\usebox\pandoc@box}%
  \else\usebox{\pandoc@box}\fi}
\def\fps@figure{htbp}
\makeatother
\providecommand{\tightlist}{\setlength{\itemsep}{0pt}\setlength{\parskip}{0pt}}

% ---- references: author-year Harvard (Cell Press / Cell Reports style) ----
%   VERIFIED vs Cell Reports author guidelines (2026-06): Cell Press uses
%   author-year in-text citations (e.g. Smith et al., 2010), NOT numbered.
%   apalike approximates Cell's Harvard list; swap to cell.bst at final
%   submission if the production team requests it.
\usepackage[round,authoryear]{natbib}
\bibliographystyle{apalike}

% ---- line numbers for peer review ----
\usepackage{lineno}\linenumbers

% ---- authors / affiliations ----
\usepackage{authblk}
\renewcommand\Authands{, }
\renewcommand\Authsep{, }
\newcommand\blfootnote[1]{\begingroup\renewcommand\thefootnote{}%
  \footnotetext{#1}\addtocounter{footnote}{-1}\endgroup}

\usepackage{hyperref}
\hypersetup{unicode=true,colorlinks=true,linkcolor=blue,citecolor=blue,urlcolor=blue,
  pdftitle={A novel lytic Aeromonas bacteriophage rescues largemouth bass from lethal Aeromonas hydrophila infection},
  pdfauthor={Liang C; Zhang Y; Xue J; Zhang H; Chen Z; Ling F; Deng L; Wang G; Ru J}}

\title{\bfseries A novel lytic \emph{Aeromonas} bacteriophage
rescues largemouth bass from lethal \emph{Aeromonas hydrophila} infection}

\author[1,2]{Changshuai Liang}
\author[1,2]{YiLin Zhang}
\author[3,4]{Jinling Xue}
\author[1,2]{Haowei Zhang}
\author[3]{Zhangsendi Chen}
\author[1,2]{Fei Ling}
\author[3,4]{Li Deng\textsuperscript{*}}
\author[1,2]{Gaoxue Wang\textsuperscript{*}}
\author[3,4]{Jinlong Ru\textsuperscript{*}}

\affil[1]{Hainan Institute of Northwest A\&F University, Sanya, Hainan Province, 572024, China}
\affil[2]{College of Animal Science and Technology, Northwest A\&F University, Xinong Road 22nd, Yangling, Shaanxi, 712100, China}
\affil[3]{Technical University of Munich; Prevention of Microbial Diseases, TUM School of Life Sciences; Central Institute of Infection Prevention (ZIP), Freising, 85354, Germany}
\affil[4]{Research unit of Microbial Diseases Prevention, Institute of Virology, Helmholtz Centre Munich -- German Research Centre for Environmental Health, Neuherberg, 85764, Germany}

\date{}

\begin{document}
\maketitle
\blfootnote{\textsuperscript{*}Corresponding authors:
  li.deng@tum.de (L.D.); wanggaoxue@126.com (G.W.); jinlong.ru@tum.de (J.R.).
  ORCID (J.R.): 0000-0002-6757-6018}

% ============================================================
%  Cell Reports front matter
% ============================================================

\section*{Highlights}
\begin{itemize}\tightlist
\item Lytic \emph{Aeromonas hydrophila} phage Hp3 forms a deeply divergent, unclassified genus
\item Its 84.7-kb genome carries no resistance or high-confidence virulence genes
\item A single phage dose raised 7-day survival from 15\% to 85\% under lethal challenge
\item Hp3 is a genetically distinct candidate for aquaculture phage biocontrol
\end{itemize}

% eTOC blurb: Cell Reports limit = 40 words. This is 39.
\section*{In Brief}
Liang et al.\ isolate Hp3, a lytic phage of \emph{Aeromonas hydrophila} that
defines a deeply divergent, unclassified viral genus. Lacking resistance and
virulence genes, Hp3 rescues largemouth bass from lethal infection, offering a
distinct candidate for aquaculture phage biocontrol.

\section*{Summary}
\noindent
Antimicrobial resistance in \emph{Aeromonas hydrophila}, a leading pathogen of
freshwater aquaculture, is eroding antibiotic control and motivating
bacteriophage-based alternatives. Here we isolate and characterize
vB\_AhyM\_Hp3 (Hp3), a lytic phage of \emph{A.~hydrophila}. Integrated
nucleotide- and proteome-level analyses place Hp3 and its only close relative
in a deeply divergent, currently unclassified \emph{Caudoviricetes} genus,
indicating that aquaculture-associated \emph{Aeromonas} phages harbour
substantial unsampled diversity. Hp3 completes a rapid (\textasciitilde90-min)
lytic cycle, remains stable across the temperature and pH ranges typical of
aquaculture, and carries an 84.75-kb genome with the hallmarks of a virulent
lifestyle but no antimicrobial-resistance genes and no high-confidence
virulence determinants. In a lethal largemouth bass (\emph{Micropterus
salmoides}) challenge, a single dose of Hp3 raised seven-day survival from
15\% to 85\% and markedly reduced liver and spleen pathology, while phage
alone caused no detectable harm. Hp3 is therefore a genetically distinct,
genomically defined candidate for aquaculture biocontrol and a tractable
starting point for mechanistic and formulation studies of phage therapy
against \emph{A.~hydrophila}.

\bigskip
% INITIAL submission: narrative Methods (body.tex) + the single Data-availability
% statement in Declarations. Cell Reports does NOT require a specific format at
% initial submission (official guidelines), so STAR Methods is deferred to revision
% and the Data-availability redundancy is avoided.
% REVISION: switch to \input{body-cr} (drops narrative Methods; loads
% star_methods.tex = Key Resources Table + Data and code availability).
    

\section{Introduction}\label{introduction}

\emph{Aeromonas hydrophila}, a Gram-negative bacterium widely
distributed in aquatic environments, represents a significant and
escalating threat to global aquaculture \citep{2023-Review_Semwal}. As
the causative agent of severe septicemic and ulcerative diseases in a
broad spectrum of hosts, it is responsible for devastating outbreaks and
substantial economic losses, particularly in the intensive farming of
economically vital species like the largemouth bass (\emph{Micropterus
salmoides})
\citep{2006-Importance_Daskalov, 2021-MiRNAseq_Zhou, 2025-Vivo_Wang}.
Conventional antibiotic control of these
infections is increasingly ineffective. Overuse of these drugs
has driven the emergence of multidrug-resistant \emph{A.
hydrophila} strains, rendering treatments ineffective, and has also led
to environmental contamination \citep{2025-Global_Jeamsripong}. Together,
antimicrobial resistance and ecological damage motivate the search for
alternatives to antibiotics.

Bacteriophage therapy is one alternative under active investigation.
Phages are natural predators of bacteria, offering a highly
specific mechanism of action: they infect and lyse their target pathogen
while leaving the host's beneficial microbiota undisturbed
\citep{2019-Phage_GordilloAltamirano}. Unlike antibiotics, phages can be
effective against biofilm-encased bacteria and are considered
environmentally safe \citep{2023-Potential_Guo}. These properties make
them well suited to controlling bacterial pathogens
in sensitive ecosystems such as aquaculture \citep{2021-Phage_Ramos-Vivas}.

However, the promise of phage therapy can only be realized through the
rigorous discovery and characterization of suitable phage candidates.
Developing an environmental isolate into a therapeutic agent is
difficult, because phages may be environmentally unstable and because
phage--bacterium co-evolution erodes efficacy. Lytic efficiency,
stability, genomic safety, and efficacy in a relevant clinical model
must all be assessed
\citep{2025-Reevaluating_Cook, 2018-Phage_Torres-Barcelo}. This study
was therefore undertaken to identify and characterize a
novel phage against \emph{A. hydrophila}. We report the isolation,
multi-level characterization, and genomic analysis of phage
vB\_AhyM\_Hp3 (Hp3), and its
therapeutic efficacy in a lethal largemouth bass infection model.

\section{Experimental Procedures}\label{experimental-procedures}

\subsection{Bacterial Strain and Phage
Isolation}\label{bacterial-strain-and-phage-isolation}

The host bacterium, \emph{Aeromonas hydrophila} strain MS01, was
isolated in our laboratory from a diseased largemouth bass
(\emph{Micropterus salmoides}). Its species identity was confirmed via
whole-genome sequencing. For host range analysis, nineteen additional
\emph{A. hydrophila} strains were used (Table~\ref{tbl-host-range}). All
strains were routinely cultured in Luria-Bertani (LB) broth at 37°C with
shaking.

Phage vB\_AhyM\_Hp3 (Hp3) was isolated from municipal sewage using an
enrichment protocol. Briefly, a logarithmic-phase culture of \emph{A.
hydrophila} MS01 (\(\mathrm{OD}_{600} \approx 0.5\)) was incubated
overnight with an equal volume of 0.22 μm-filtered sewage. The culture
was then centrifuged (4000 × g, 20 min, 4°C), and the supernatant was
filter-sterilized. A single phage plaque was isolated and purified
through three successive rounds of the double-layer agar method to
ensure clonal purity. High-titer phage stocks were prepared and stored
at 4°C for subsequent experiments.

\subsection{Transmission Electron Microscopy
Analysis}\label{transmission-electron-microscopy-analysis}

To visualize phage morphology, a high-titer phage suspension (approx.
10\textsuperscript{9} PFU/mL) was adsorbed onto a carbon-coated copper
grid. The grid was then negatively stained with 3\% (w/v) uranyl acetate
and allowed to air-dry. Phage particles were examined using a Hitachi
HT7800 transmission electron microscope to determine their
ultrastructural characteristics.

\subsection{Phage Host Range
Determination}\label{phage-host-range-determination}

To evaluate the lytic activity of phage Hp3 against pathogenic bacterial
strains in aquaculture, nineteen \emph{A. hydrophila} strains isolated
from diseased fish in different geographical locations
(Table~\ref{tbl-host-range}) were retrieved from storage at -80°C and
cultured in LB liquid medium at 37°C with shaking at 180 rpm. Following
bacterial resuscitation, the host range of phage Hp3 was determined
using standard spot assay.

\subsection{Determination of optimal Multiplicity of Infection
(MOI)}\label{determination-of-optimal-multiplicity-of-infection-moi}

A suspension of \emph{A. hydrophila} in the late logarithmic growth
phase was prepared as the test concentration. Phage Hp3 was added to the
bacterial suspension at MOI of 0.001, 0.01, 0.1, 1.0, 10, 100, and 1000,
and then incubated at 37°C for 3 hours. The co-culture was then
centrifuged at 4000 × g for 20 min, and the supernatant was filtered
through a 0.22 μm membrane to remove bacteria. The filtrate was
incubated with the \emph{A. hydrophila} culture, and the phage titer was
determined using the double-layer agar plate method. The MOI that
produced the highest phage titer was considered the optimal MOI, which
was used in subsequent experiments.

\subsection{One-Step Growth Curve}\label{one-step-growth-curve}

The one-step growth curve was performed to determine the latent period
of Hp3. Phage Hp3 was added to the \emph{A. hydrophila} suspension and
incubated at 37°C for 10 min. The mixture was then centrifuged at 4000 ×
g for 20 min, and the supernatant was discarded. The bacterial pellet
was resuspended in fresh medium and incubated at 37°C with shaking at
180 rpm. Samples (200 μL) were collected at different time points, and
the phage titer was determined using the double-layer agar plate method.

\subsection{Environmental Stability
Assays}\label{environmental-stability-assays}

To assess thermal stability, phage stocks were incubated at various
temperatures (4, 28, 37, 45, 60, and 80°C) for one hour. For pH
stability, phages were incubated for one hour at 37°C in an SM buffer
adjusted to pH values ranging from 4.0 to 12.0. Following incubation,
the remaining viable phage titer in each sample was determined using the
double-layer agar method. The experiment was repeated three times.

\subsection{\texorpdfstring{\emph{In Vitro} Lysis
Assay}{In Vitro Lysis Assay}}\label{in-vitro-lysis-assay}

The lytic capability of Hp3 was monitored in a 96-well plate format.
MS01 cultures were infected with Hp3 at different MOIs (0.01 to 100).
The optical density at 600 nm (\(\mathrm{OD}_{600}\)) of the cultures
was recorded every 20 minutes for 48 hours using an automated microplate
reader. Each MOI was assayed in three independent biological replicates and the
uninfected negative control in six.

\subsection{Whole Genome Analysis}\label{whole-genome-analysis}

Phage DNA was extracted from the enriched phages using Norgen phage DNA
kit (cat. 46800) following manufacturer instructions. The concentration
and purity of DNA were determined using a Nanodrop spectrophotometer and
1\% agarose gel electrophoresis. The extracted phage DNA was sequenced
on Illumina NovaSeq 6000 platform. Raw sequencing reads were
quality-controlled using fastp \citep{2018-Fastp_Chen}, and then
assembled using Unicycler v0.5.1
(\url{https://github.com/rrwick/Unicycler}) to obtain the complete
genome sequence.

The assembled complete genome was annotated using Pharokka v1.3.2
\citep{2022-Pharokka_Bouras}, Phold v0.2 \citep{2026-Protein_Bouras} and
Phynteny v0.1.13 \citep{2024-Susiegriggo_SusieGrigson}. For taxonomic
classification, a proteomic network was constructed using vConTACT3
(\url{https://bitbucket.org/MAVERICLab/vcontact3}) against the NCBI
viral RefSeq database (v230). The resulting network, which connects
genomes sharing at least one homologous protein cluster, was visualized
in Cytoscape \citep{2003-Cytoscape_Shannon}; a subgraph showing only
viruses connected to Hp3 is presented (Supplementary Figure S1).
Edge thickness in the network is proportional to the number of shared
protein clusters. Genus- and species-level relationships were quantified
from whole-genome intergenomic similarities computed with the VIRIDIC
algorithm \citep{2020-VIRIDIC_Moraru} (genus and species demarcation
thresholds of 70\% and 95\%, respectively). Genome-wide relatedness of
Hp3 to all complete \emph{Aeromonas} phage genomes in GenBank (n = 200)
was screened with Mash \citep{2016-Mash_Ondov}. To place Hp3 within the
broader double-stranded DNA virosphere, a whole-proteome tree against
5,633 reference prokaryotic dsDNA viruses was constructed with ViPTree
\citep{2017-ViPTree_Nishimura}, which scores pairwise relatedness as a
normalized tBLASTx genomic similarity (\(S_G\)). Genome-based phylogeny
and genus-level OPTSIL boundaries were further assessed with VICTOR
\citep{2017-VICTOR_Meier-Kolthoff}, applying the Genome-BLAST Distance
Phylogeny (GBDP) method to amino-acid sequences (recommended formula D6)
for Hp3, HJ05, their nearest unclassified ViPTree neighbours, and the
type species of eight major \emph{Caudoviricetes} families. Higher-rank
OPTSIL partitions were inspected but not used as family-level evidence
because family-level resolution is known to be weak for distant
prokaryotic viruses in VICTOR and, in our calibration panel, merged
several ICTV-recognized families.

Comparative
genome map was generated using clinker \citep{2021-Clinker_Gilchrist},
displaying homologous gene pairs sharing \textgreater30\% identity at
the amino acid level. Potential antibiotic-resistance and virulence
genes were screened against the CARD \citep{2020-CARD_Alcock} and VFDB
\citep{2019-VFDB_Liu} databases, respectively. The replication lifestyle
(virulent versus temperate) was predicted with BACPHLIP
\citep{2021-BACPHLIP_Hockenberry}, which infers lifestyle from a set of
conserved lifestyle-associated protein domains, and was cross-checked by
manual inspection of the genome annotation for lysogeny-associated genes
(a site-specific integrase and a lysogenic-maintenance repressor).

\subsection{\texorpdfstring{\emph{In Vivo} Therapeutic Efficacy in
Largemouth
Bass}{In Vivo Therapeutic Efficacy in Largemouth Bass}}\label{in-vivo-therapeutic-efficacy-in-largemouth-bass}

Healthy largemouth bass (\emph{Micropterus salmoides}) were purchased
from a commercial fish farm in Guangzhou (Guangdong Province, China).
Upon arrival, the fish (average length 5.62 ± 0.52 cm; average weight
2.38 ± 0.31 g) were allowed to acclimatize to laboratory conditions for
two weeks. During this period, they were housed in aerated tanks with a
recirculating water system at 25 °C, pH of 6.8-8.0, and dissolved oxygen
\textgreater5.7 mg/L and fed a commercial diet (Maoming Hailong Feed
Co., Ltd.) twice daily. Twenty healthy largemouth bass in each group
were intraperitoneally injected with 0.1 mL suspension of \emph{A.
hydrophila} at various concentrations of 1 × 10\textsuperscript{6}, 1 ×
10\textsuperscript{7}, 1 × 10\textsuperscript{8}, 1 ×
10\textsuperscript{9}, and 1 × 10\textsuperscript{10} CFU/mL. The fish
were observed for seven consecutive days, and the cumulative mortality
rate was recorded. The LD\textsubscript{50} calculated using the
modified Käber method was 4.47 × 10\textsuperscript{7} CFU/mL. Based on
this, a bacterial concentration of 9 × 10\textsuperscript{7} CFU/mL was
selected as the challenge dose for the formal experiment.

Eighty healthy largemouth bass were randomly divided into four groups of
20 each: the phage group, the bacterial group, the treatment group, and
the PBS control group. In the phage group, 0.1 mL of PBS was injected
intraperitoneally, followed by 0.1 mL of a phage Hp3 suspension (9 ×
10\textsuperscript{8} PFU/mL) 1 hour later. In the bacterial group, 0.1
mL of \emph{A. hydrophila} suspension (9 × 10\textsuperscript{7} CFU/mL)
was injected intraperitoneally, followed by 0.1 mL of PBS 1 hour later.
In the PBS group, 0.1 mL of PBS was injected intraperitoneally at 0 and
1 hour, respectively. In the treatment group, 0.1 mL of \emph{A.
hydrophila} suspension (9 × 10\textsuperscript{7} CFU/mL) was injected
intraperitoneally, followed by 0.1 mL of a phage Hp3 suspension (9 ×
10\textsuperscript{7} PFU/mL) one hour later.

After injection, fish in each group were housed in isolated water
maintained at 25 ± 1°C. Daily survival was recorded for 7 consecutive
days, and the cumulative survival rate was calculated to evaluate the
\emph{in vivo} therapeutic effect of the phage. Three days after
infection, tissue sampling was performed. Prior to sampling, fish were
anesthetized using tricaine methanesulfonate (MS-222) (Shanghai Aladdin
Bio-Chem Technology Co., Ltd., Shanghai, China) at a concentration of 50
mg/L to minimize stress. For tissue collection, fish were euthanized by
immersion in an overdose of MS-222 at 350 mg/L, followed by verification
of death (cessation of opercular movement). The liver and spleen tissues
of fish in each group were then collected, fixed in 4\%
paraformaldehyde, embedded in conventional paraffin, and sliced (3 μm)
for histopathological observation under an optical microscope after
hematoxylin-eosin staining. All procedures were performed in strict
accordance with the ethical protocols approved by Northwest A\&F
University and adhered to the American Veterinary Medical Association
(AVMA) Guidelines for the Euthanasia of Animals (2020). The \emph{in
vivo} study is reported with reference to the ARRIVE guidelines (Animal
Research: Reporting of In Vivo Experiments), with unresolved reporting
limitations for power calculation, allocation procedures, and blinding
described in the Discussion.

\subsection{Statistical analysis}\label{statistical-analysis}

Unless otherwise specified in the figure legends, all quantitative data
from experiments with multiple biological replicates are presented as the mean ±
standard deviation (SD). The Kaplan-Meier method was used to plot
survival curves. Survival differences were first assessed using an
overall k-sample Log-rank (Mantel-Cox) test across the PBS, phage-only,
bacterial, and treatment groups. Pairwise Log-rank comparisons were then
performed with Benjamini-Hochberg adjustment for multiple testing. The
treatment effect was additionally estimated using a Cox
proportional-hazards model fitted to the bacterial and treatment groups
(with the bacterial group as the reference); the phage-only and PBS
groups were excluded from this model as they recorded no deaths. For bacterial growth (in
vitro lysis) assays, differences in optical density between groups at
the experimental endpoint (48 h) were analyzed using one-way analysis of
variance (ANOVA) followed by Tukey's honestly significant difference (HSD)
post-hoc test for pairwise comparisons. A p-value \textless{} 0.05 was
considered statistically significant. All statistical analyses were performed using
R software (v4.5.1).

\section{Results}\label{results}

\subsection{Purification and Morphological Characterization of Phage
Hp3}\label{purification-and-morphological-characterization-of-phage-hp3}

Hp3 forms single round, small, and transparent plaques with no lytic
halo around them on a double-layer agar plate
(Figure~\ref{fig-hp3-morphology} A). TEM revealed that Hp3 possesses an
icosahedral head and a contractile tail, a morphology characteristic of
the class \emph{Caudoviricetes}. The head diameter is about 81 nm, the
tail length is about 143 nm, and the width is about 23 nm.
(Figure~\ref{fig-hp3-morphology} B)

\begin{figure}

\centering{

\pandocbounded{\includegraphics[keepaspectratio]{../analyses/data/121-figs/fig1.png}}

}

\caption{\label{fig-hp3-morphology}Morphological observations of phage
Hp3. (A) The plaque morphology of bacteriophage Hp3 observed on
double-layer agar plates; (B) The TEM of bacteriophage Hp3 (Bar = 100
nm).}

\end{figure}%

\subsection{Hp3 Has a Narrow Host
Range}\label{hp3-exhibits-potent-lytic-activity-and-a-narrow-host-range}

The lytic spectrum of Hp3 was determined using 19 \emph{A. hydrophila}
strains (Table~\ref{tbl-host-range}). Hp3 exhibited a narrow lytic
range, successfully lysing only two bacterial strains.

\begin{longtable}[]{@{}
  >{\raggedright\arraybackslash}p{(\linewidth - 10\tabcolsep) * \real{0.2933}}
  >{\raggedright\arraybackslash}p{(\linewidth - 10\tabcolsep) * \real{0.1067}}
  >{\centering\arraybackslash}p{(\linewidth - 10\tabcolsep) * \real{0.0933}}
  >{\raggedright\arraybackslash}p{(\linewidth - 10\tabcolsep) * \real{0.3067}}
  >{\raggedright\arraybackslash}p{(\linewidth - 10\tabcolsep) * \real{0.1067}}
  >{\centering\arraybackslash}p{(\linewidth - 10\tabcolsep) * \real{0.0933}}@{}}
\caption{Host range of phage Hp3 against \emph{Aeromonas hydrophila}
strains from different geographical regions. +, susceptible to lysis; −,
resistant to lysis.}\label{tbl-host-range}\tabularnewline
\toprule\noalign{}
\begin{minipage}[b]{\linewidth}\raggedright
Bacteria strain
\end{minipage} & \begin{minipage}[b]{\linewidth}\raggedright
Origin
\end{minipage} & \begin{minipage}[b]{\linewidth}\centering
Lysis
\end{minipage} & \begin{minipage}[b]{\linewidth}\raggedright
Bacteria strain
\end{minipage} & \begin{minipage}[b]{\linewidth}\raggedright
Origin
\end{minipage} & \begin{minipage}[b]{\linewidth}\centering
Lysis
\end{minipage} \\
\midrule\noalign{}
\endfirsthead
\toprule\noalign{}
\begin{minipage}[b]{\linewidth}\raggedright
Bacteria strain
\end{minipage} & \begin{minipage}[b]{\linewidth}\raggedright
Origin
\end{minipage} & \begin{minipage}[b]{\linewidth}\centering
Lysis
\end{minipage} & \begin{minipage}[b]{\linewidth}\raggedright
Bacteria strain
\end{minipage} & \begin{minipage}[b]{\linewidth}\raggedright
Origin
\end{minipage} & \begin{minipage}[b]{\linewidth}\centering
Lysis
\end{minipage} \\
\midrule\noalign{}
\endhead
\bottomrule\noalign{}
\endlastfoot
\emph{A. hydrophila} N8-2 & Liaoning & + & \emph{A. hydrophila} 8-2 &
Chongqing & − \\
\emph{A. hydrophila} 6-1 & Chongqing & + & \emph{A. hydrophila} 8-6 &
Chongqing & − \\
\emph{A. hydrophila} N2-2 & Liaoning & − & \emph{A. hydrophila} F03-1 &
Shaanxi & − \\
\emph{A. hydrophila} N2-3 & Liaoning & − & \emph{A. hydrophila} F05-5 &
Shaanxi & − \\
\emph{A. hydrophila} N3-2 & Liaoning & − & \emph{A. hydrophila} F05-7 &
Shaanxi & − \\
\emph{A. hydrophila} N3-3 & Liaoning & − & \emph{A. hydrophila} D6-1 &
Guangdong & − \\
\emph{A. hydrophila} N4-3 & Liaoning & − & \emph{A. hydrophila} D2-5 &
Guangdong & − \\
\emph{A. hydrophila} N7-5 & Liaoning & − & \emph{A. hydrophila} NS9-2 &
Jiangsu & − \\
\emph{A. hydrophila} N7-6 & Liaoning & − & \emph{A. hydrophila} NS10-1 &
Jiangsu & − \\
\emph{A. hydrophila} GS9-4 & Guizhou & − & & & \\
\end{longtable}

\subsection{Hp3 Has a Rapid Lytic Cycle and Is Stable Across Temperature
and pH}\label{hp3-possesses-a-rapid-lytic-cycle-and-broad-environmental-stability}

\subsubsection{Optimal MOI}\label{optimal-moi}

The multiplicity of infection (MOI) refers to the ratio of phages that
can be adsorbed to host bacteria within a specific period of time, and
is an important indicator of phage infection experiments. In order to
study the optimal multiplicity of infection of phages, this experiment
set 7 MOIs (0.001, 0.01, 0.1, 1.0, 10, 100, 1000). The results of the
optimal phage MOI determination showed that when MOI=1.0, the titer of
Hp3 reached the maximum (Figure~\ref{fig-hp3-properties} A).

\subsubsection{One-step growth curve of
phage}\label{one-step-growth-curve-of-phage}

The one-step growth curve revealed a latent period of approximately 45
minutes, followed by a burst phase of about 45 minutes, resulting in a
complete lytic cycle of roughly 90 minutes.
(Figure~\ref{fig-hp3-properties} B).

\subsubsection{Phage stability analysis}\label{phage-stability-analysis}

To assess its suitability for application, the stability of phage Hp3
was evaluated under various thermal and pH conditions. Hp3 showed
high thermal stability, maintaining a consistently high titer of
approximately 10\textsuperscript{8} PFU/mL after a one-hour incubation
at temperatures ranging from 4°C to 45°C. The phage's viability began to
decline at 60°C, where the titer dropped by two orders of magnitude, and
it was completely inactivated at 80°C (Figure~\ref{fig-hp3-properties}
C).

Hp3 was also stable across a wide pH range. The
phage titer remained stable at approximately 10\textsuperscript{8}
PFU/mL after one hour of incubation in buffers with pH values from 6.0
to 9.0. A moderate, one-log reduction in titer was observed at pH 10.0,
while the phage was completely inactivated at pH values below 5.0 and
above 11.0 (Figure~\ref{fig-hp3-properties} D). These results indicate
that Hp3 is tolerant to the range of temperature and pH conditions
commonly found in aquaculture environments.

\begin{figure}

\centering{

\pandocbounded{\includegraphics[keepaspectratio]{../analyses/data/121-figs/fig2.png}}

}

\caption{\label{fig-hp3-properties}Biological properties of Hp3. (A)
Determination of optimal infection multiplicity of phage Hp3. (B)
One-step growth curve of phage Hp3 on \emph{A. hydrophila} at 37°C. The
figure shows the phage titer at different time points. (C) Temperature
stability of Hp3. (D) pH stability of Hp3. The experiment was repeated
three times. The data are presented as the mean ± SD of three
replicates.}

\end{figure}%

\subsubsection{Phage lysis curve}\label{phage-lysis-curve}

To evaluate the lytic efficacy of Hp3 over time, we monitored its effect
on the growth of \emph{A. hydrophila} at various MOIs for 48 hours. As
shown in Figure~\ref{fig-lysis-curve}, Hp3 demonstrated a significant
dose-dependent antibacterial effect, maintaining the bacterial density
below that of the untreated control throughout the experiment. At the
48-hour endpoint, the optical density in the different treatment groups
showed statistically significant differences (One-way ANOVA, F(5, 15) =
23.191, p \textless{} 0.001). All five MOIs suppressed bacterial density
relative to the untreated control (Tukey HSD, all p \textless{} 0.005), with
MOI 1.0 and 10 producing the strongest and most sustained suppression (not
significantly different from each other, p \textgreater{} 0.99, and
significantly stronger than MOI 100, p \textless{} 0.05). In these
conditions, an initial phase of bacterial growth was followed by a sharp
decline in density, although regrowth was observed at later time points,
likely due to the emergence of phage-resistant bacteria. Despite this
regrowth, Hp3 suppressed \emph{A. hydrophila} across the tested MOIs,
supporting its lytic activity against the strain.

\begin{figure}

\centering{

\pandocbounded{\includegraphics[keepaspectratio]{../analyses/data/121-figs/fig3.png}}

}

\caption{\label{fig-lysis-curve}48h lysis curve of Hp3. five MOI values
(0.01, 0.1, 1.0, 10 and 100) were incubated with the control group at
37°C. The phage titer was determined by measuring the
\(\mathrm{OD}_{600}\) absorbance. The data are presented as the mean ± SD (n = 3 biological replicates per MOI group; n = 6 for the
untreated control).
Statistical analysis indicated significant differences in optical
density between groups at the 48-hour endpoint (One-way ANOVA, F(5, 15)
= 23.191, p \textless{} 0.001)}

\end{figure}%

\subsection{Genomic Analysis Places Hp3 in a Deeply Divergent
\emph{Caudoviricetes}
Genus}\label{genomic-analysis-reveals-hp3-represents-a-putative-novel-viral-family}

Whole-genome sequencing revealed that phage Hp3 possesses a
double-stranded DNA genome of 84,750 bp (GenBank accession PX754746.1)
with a GC content of 59\% (Figure~\ref{fig-genomics} A). Bioinformatic analysis predicted a total
of 105 open reading frames (ORFs), of which 89 were assigned putative
functions, while 16 were annotated as hypothetical proteins.

To determine the taxonomic position of Hp3, we integrated
nucleotide-level and whole-proteome analyses. At the
nucleotide level, VIRIDIC placed Hp3 and \emph{Aeromonas} phage HJ05
(GenBank accession OQ594355) within a single genus, sharing 83.65\%
intergenomic similarity---above the 70\% genus threshold but below the
95\% species threshold---indicating that the two are distinct species of
the same genus. Across all 200 complete
\emph{Aeromonas} phage genomes in GenBank, Hp3 showed appreciable
similarity to HJ05 alone (Mash distance = 0.12), whereas the Mash
distance to every other \emph{Aeromonas} phage was 1.0 (no shared
k-mers), underscoring the isolation of this genus. To test whether the
genus falls within any established family, we placed Hp3 in a
whole-proteome tree of 5,633 reference prokaryotic dsDNA viruses
(ViPTree): Hp3 branched as an isolated, deep lineage, its highest
genomic similarity to any reference being only \(S_G\) = 0.073 (to the
unclassified \emph{Rhizobium} phage vB\_RleS\_L338C), with no member of
any recognized family approaching it. Consistently, the vConTACT3
protein-sharing network grouped Hp3 with HJ05 and, more distantly,
vB\_RleS\_L338C, a cluster clearly separated from all established
families including its nearest neighbors in the \emph{Peduoviridae} and
\emph{Straboviridae} (Supplementary Figure S1); we note, however,
that L338C shares only \textasciitilde3\% genome-wide nucleotide
identity with Hp3 (VIRIDIC), indicating that this protein-sharing link
reflects a small number of conserved proteins rather than close kinship.
As HJ05 is itself an unclassified member of the \emph{Caudoviricetes},
these convergent analyses indicate that Hp3 and HJ05 constitute a new,
deeply divergent genus. This genus-level assignment was independently
corroborated by amino-acid genome-BLAST distance phylogeny
(VICTOR/GBDP, formula D6; Figure~\ref{fig-genomics} B), in which Hp3 and HJ05 again formed an
exclusive genus-level cluster (pseudo-bootstrap = 100), well separated
from the type species of eight major \emph{Caudoviricetes} families. We
therefore did not use VICTOR family-level OPTSIL partitions as support
for family rank. In this saturated-divergence panel, that partition
merged \emph{Straboviridae}, \emph{Ackermannviridae},
\emph{Demerecviridae}, and \emph{Herelleviridae}, illustrating the
known weak family-level resolution of VICTOR for distant prokaryotic
viruses \citep{2017-VICTOR_Meier-Kolthoff}. Because family-level
demarcation in \emph{Caudoviricetes} has no universal quantitative
threshold and formal family creation generally requires multiple genera
plus core-gene and phylogenetic criteria
\citep{2021-Roadmap_Turner,2023-Abolishment_Turner}, we interpret the
ViPTree and vConTACT3 evidence as supporting a putative
family-level lineage pending broader sampling and formal ICTV
evaluation, rather than an established family.

A direct comparative analysis of phage Hp3 with HJ05 revealed a high
degree of genomic identity and gene order conservation (synteny) across
the majority of their genomes, confirming their close evolutionary
relationship (Figure~\ref{fig-genomics} C). A notable difference exists
at the 5' end of the genome, where Hp3 carries a block of additional
genes absent in HJ05 (Figure~\ref{fig-genomics} A, dashed rectangle);
the majority of these encode hypothetical proteins of unknown function,
and we therefore refrain from assigning them a specific role.

Beyond the structural modules typical of the \emph{Caudoviricetes}, the
Hp3 genome carries nucleotide-metabolism genes that may support phage
DNA replication, including a thymidylate synthase (hp3\_0036,
\emph{thyA}) and an anaerobic ribonucleotide reductase (hp3\_0058,
\emph{nrdD}). These genes are common replication-associated features of
lytic phages rather than evidence for broad host-metabolism remodeling.
A less common annotation is hp3\_0002, a putative, distantly related
DAHP-synthase structural homolog (annotated as
phospho-2-dehydro-3-deoxyheptonate aldolase; Foldseek structural
homology to characterized DAHP synthases, fident = 0.38, E-value = 5 ×
10\textsuperscript{-36}). hp3\_0002 lies within the complete circular
Hp3 genome, adjacent to the terminase large-subunit gene and followed by
additional phage-annotated ORFs, rather than at an unresolved contig
edge. These assignments rest on sequence and structural homology and
await experimental validation; the enzymatic activity, substrate
specificity, and physiological role of hp3\_0002 during infection remain
unknown.

Consistent with the rapid lysis observed \emph{in vitro}, the Hp3 genome
carries the hallmarks of a virulent (lytic) rather than a temperate
lifestyle. It encodes a complete lysis cassette---two holins (hp3\_0069 and
hp3\_0099) and an amidase-family endolysin (hp3\_0070)---alongside a
T4-like homologous-recombination and replication system, including a
UvsX-like recombinase (hp3\_0091) and RuvC-like Holliday-junction
resolvases. Conversely, it lacks the core machinery required to establish
and maintain lysogeny: no tyrosine- or serine-family site-specific
integrase and no CI-like lysogenic-maintenance repressor or immunity
region. The single annotation falling in the integration-and-excision
category is one open reading frame at the extreme 3' terminus of the
genome (hp3\_0105), assigned as a transposase only by low-confidence remote
structural homology (Phold/Foldseek structural identity = 0.24, with no
significant sequence-homology hit from Pharokka) and most parsimoniously
interpreted as a horizontally acquired mobile-element remnant rather than a
functional prophage-integration system; a solitary transposase-like gene
does not confer lysogenic capacity. Concordantly, BACPHLIP classified Hp3
as virulent with high confidence (lytic probability = 0.97)
\citep{2021-BACPHLIP_Hockenberry}. We therefore designate Hp3 a lytic phage
on the basis of converging genomic and phenotypic evidence.

Finally, a genomic safety screen detected no antibiotic-resistance genes
(CARD). The only virulence-database flag, hp3\_0022, arose solely from
distant structural homology---a Phold/Foldseek match to VFDB entry
VFG047113 at 16\% amino-acid identity (fident = 0.16)---with no
corresponding sequence-homology hit in the Pharokka VFDB screen, and is
therefore unlikely to represent a functional virulence determinant.
Together with the absence of lysogeny machinery described above, this
genomic-safety profile supports Hp3 as a candidate for therapeutic
evaluation, while uncharacterized proteins, including hp3\_0002 and
hp3\_0022, remain for future functional validation.

\begin{figure}

\centering{

\pandocbounded{\includegraphics[keepaspectratio]{../analyses/data/121-figs/fig4.png}}

}

\caption{\label{fig-genomics}Comprehensive genomic, taxonomic, and
comparative analyses of phage Hp3. (A) Circular map of the Hp3 genome.
Functional gene categories are based on the PHROGS database
\citep{2021-PHROG_Terzian}. Rings from outside in represent: ORFs, GC
content, and GC skew. (B) Amino-acid genome-BLAST distance phylogeny
(VICTOR, GBDP formula D6) of Hp3, HJ05, their nearest ViPTree
neighbours, and the type species of eight major \emph{Caudoviricetes}
families (\emph{Peduoviridae}, P2; \emph{Ackermannviridae}, CBA120;
\emph{Straboviridae}, T4; \emph{Demerecviridae}, T5; \emph{Herelleviridae},
SPO1; \emph{Autographiviridae}, T7; \emph{Schitoviridae}, N4; and
\emph{Drexlerviridae}, T1; for compactness the \emph{Ackermannviridae}
representative is shown by strain name only, CBA120, in the panel). Node
values are pseudo-bootstrap support, the coloured columns show OPTSIL
clusters at the family, subfamily, genus, and species ranks (left to
right), and the grey bars at the right give the total amino-acid sequence
length of each genome's proteome (range 10,111--52,048 aa). Hp3 and HJ05 form an exclusive, fully supported genus
(pseudo-bootstrap = 100). The family-rank OPTSIL partition is shown for
completeness but is unreliable in this saturated-divergence panel---it
merges the recognized families \emph{Straboviridae},
\emph{Ackermannviridae}, \emph{Demerecviridae}, and
\emph{Herelleviridae}---and was therefore not used to infer family rank.
(C) Linear genome comparison of Hp3 and its closest relative,
\emph{Aeromonas} phage HJ05. Shaded regions connect homologous genes
with \textgreater30\% amino acid sequence identity. The dashed
rectangles in (A) and (C) highlight a region of genomic difference
compared with HJ05.}

\end{figure}%

\subsection{Phage Hp3 Rescues Largemouth Bass from a Lethal A.
hydrophila
Challenge}\label{phage-hp3-rescues-largemouth-bass-from-a-lethal-a.-hydrophila-challenge}

The in vivo experimental design is illustrated in
Figure~\ref{fig-therapy} A, and the resulting survival curves for each
group are shown in Figure~\ref{fig-therapy} B. In the phage-only group,
intraperitoneal injection of a high dose of Hp3 resulted in 100\%
survival over 7 days, indicating no acute mortality under these dosing
and timeframe conditions. In contrast, the group infected with
\emph{A. hydrophila} began to exhibit mortality within 24 hours
post-infection, and by day three, the survival rate had dropped to
15\%, where it remained for the rest of the experiment. In the treatment
group, however, administration of Hp3 one hour after infection resulted
in a final survival rate of 85\%. Survival differed strongly across the
four groups (overall Log-rank test, \(\chi^2\) = 64.6, df = 3,
p \textless{} 0.001; exact \(p = 6.2 \times 10^{-14}\)). In
Benjamini-Hochberg-adjusted pairwise comparisons, treatment significantly
improved survival relative to the bacterial group (Log-rank test,
\(\chi^2\) = 20.474, \(p_{\mathrm{adj}}\) = 1.2 ×
10\textsuperscript{-5}), whereas treatment did not differ significantly
from the PBS or phage-only controls (both \(p_{\mathrm{adj}}\) = 0.09);
the PBS and phage-only curves were identical. A Cox proportional-hazards
model using the bacterial group as the reference estimated a strong
reduction in mortality hazard for treated fish (HR = 0.091, 95\% CI =
0.026--0.320, \(p = 1.8 \times 10^{-4}\)).

Histopathological observations (Figure~\ref{fig-therapy} C-F) were
consistent with the survival results. The liver and spleen structures of
the PBS and phage control groups were intact, with no obvious lesions.
The liver tissue of the infected group showed diffuse vacuolar
degeneration, with swollen hepatocytes, nuclear displacement, blurred
nuclear and cytoplasmic structures and cell boundaries, and irregular
cell arrangement. Inflammatory cell infiltration was observed in the
spleen tissue, with scattered necrotic foci in both the red and white
pulp, and unclear boundaries of the parenchymal cells. In contrast, the
tissue structure of the fish treated with Hp3 remained largely intact,
with only a few focal vacuolar changes, which were markedly milder
than those in the infected group. These histopathological findings
parallel the survival data: Hp3 both improved survival after infection
with \emph{A}. \emph{hydrophila} and reduced liver and spleen damage
in vivo.

\begin{figure}

\centering{

\pandocbounded{\includegraphics[keepaspectratio]{../analyses/data/121-figs/fig5.png}}

}

\caption{\label{fig-therapy}Therapeutic efficacy and histopathological
analysis of phage Hp3 in a largemouth bass infection model. (A)
Schematic of the \emph{in vivo} experimental design, illustrating the
treatment protocols for the phage-only, bacteria-only, phage-treatment,
and PBS control group. (B) Kaplan-Meier survival curves of largemouth
bass over a seven-day period following challenge with \emph{A.
hydrophila} (9 × 10\textsuperscript{7} CFU/mL). The treatment group
received a single dose of phage Hp3 (MOI = 1.0) one hour post-infection.
Survival differed across groups (overall Log-rank test, \(\chi^2\) =
64.6, df = 3, p \textless{} 0.001), and treatment improved survival
relative to the bacterial group after Benjamini-Hochberg correction
(pairwise Log-rank test, \(\chi^2\) = 20.474,
\(p_{\mathrm{adj}}\) = 1.2 × 10\textsuperscript{-5}). (C-F) Representative histopathological images (H\&E
staining) of liver (top row, C1-F1) and spleen (bottom row, C2-F2)
tissues collected three days post-infection. (C1, C2) The PBS control
group, showing normal, healthy tissue architecture. (D1, D2) The
\emph{A. hydrophila}-infected group, exhibiting severe pathology,
including diffuse vacuolar degeneration of hepatocytes (D1, red arrow),
nuclear eccentricity (D1, blue arrow), extensive inflammatory
infiltrates (D2, red arrow), and necrotic foci (D2, blue arrow). (E1,
E2) The treatment group, which showed significantly mitigated pathology
with only mild and focal tissue changes. (F1, F2) The phage-only control
group, which also displayed intact tissue structure comparable to the
PBS control. Scale bars = 100 µm.}

\end{figure}%

\section{Discussion}\label{discussion}

In this study, we isolated and systematically characterized a lytic
\emph{Aeromonas hydrophila} phage, Hp3. Beyond its therapeutic
potential, our genomic analysis supports the placement of Hp3 and its
closest relative, \emph{Aeromonas} phage HJ05, in a deeply divergent,
currently unclassified \emph{Aeromonas}-phage genus within the
\emph{Caudoviricetes}. This placement should be read in the context of
the current post-\emph{Caudovirales} taxonomy, in which many
\emph{Caudoviricetes} genera remain unclassified at family rank and
\emph{Aeromonas} phages are known to occupy divergent protein-sharing
groups \citep{2023-Abolishment_Turner,2022-Characterization_Bujak}. Hp3
also possesses several attributes favorable for further evaluation as an
antibiotic alternative in aquaculture, including a lytic cycle of
approximately 90 minutes, stability across temperature and pH ranges
relevant to aquaculture, and a genomic profile consistent with a lytic
lifestyle---lacking lysogeny-associated integrase and repressor genes and
classified as virulent by BACPHLIP---with no known
antimicrobial-resistance genes or high-confidence virulence
determinants, providing a preliminary genomic safety basis for
therapeutic assessment.

However, the long-term efficacy of any phage therapy is fundamentally
dictated by the co-evolutionary ``arms race'' between the phage and its
bacterial host \citep{2023-Developing_Oromi-Bosch}. Two of our
observations reflect this pressure. First, the narrow host range of Hp3
reflects the highly specific interaction between phage receptor-binding
proteins and bacterial surface molecules, which are the very structures
bacteria can modify to evade predation. Second, the regrowth of the
bacterial population observed in our in vitro lysis assays is consistent
with the rapid emergence of resistant variants under phage selection
pressure. The durable in vivo rescue despite this in vitro regrowth may
reflect immunophage synergy, in which phage-mediated reduction of
pathogen density allows host immunity to clear residual bacteria before
resistance becomes clinically dominant \citep{2018-Phage_Torres-Barcelo}.
Because tissue bacterial loads were not quantified as CFU/g, bacterial
clearance remains an inferred mechanism rather than a directly measured
endpoint in this study.

Hp3 should therefore not be treated as a static monophage antimicrobial
agent. Combining phages that recognize different receptors or use
distinct infection mechanisms can reduce the probability of simultaneous
resistance and improve host-range breadth \citep{2021-Phage_Abedon,
2023-Developing_Oromi-Bosch}. Hp3 may contribute to such future
multi-phage formulations, but cocktail compatibility, resistance
management, and scalable delivery remain to be tested. The present
efficacy experiment used intraperitoneal injection; field deployment in
aquaculture will require evaluation by practical delivery routes such as
immersion or feed administration \citep{2021-Phage_Ramos-Vivas}.

A deeper layer of the arms race unfolds intracellularly, where bacteria
deploy defense systems such as restriction-modification and CRISPR-Cas
to neutralize invading genomes \citep{2010-Bacteriophage_Labrie},
against which phages have evolved diverse anti-defense systems
\citep{2024-Growing_Murtazalieva}. Our in silico screen did not detect
any known anti-defense genes in the Hp3 genome, and the HJ05-absent 5'
gene block consists largely of hypothetical proteins whose functions
remain undetermined; whether any of these contribute to counter-defense
is an open question that will require direct functional testing. The
nucleotide-metabolism genes identified in Hp3 are consistent with
replication support in lytic phages. By contrast, the hp3\_0002
DAHP-synthase-like hit is a more unusual but still speculative genomic
observation: given the distant structural homology and the absence of
transcriptomic, biochemical, or infection-stage functional data, it
should not yet be interpreted as evidence for shikimate-pathway
modulation during infection.

Despite these long-term evolutionary considerations, Hp3 showed clear
therapeutic activity in the largemouth bass infection model. A single
administration of the phage significantly improved survival after lethal
\emph{A. hydrophila} challenge and was accompanied by milder liver and
spleen pathology. This result is consistent with, and builds on, an
established aquaculture phage-therapy paradigm rather than representing
an unprecedented proof of principle. For example, phage AhFM11, a
168-kb \emph{Straboviridae} phage active against hypervirulent
\emph{A. hydrophila}, achieved high protection in fish by injection,
immersion, and oral-feed delivery routes \citep{2024-Novel_MuliyaSankappa};
aquaculture phage therapy more broadly has also been reviewed as a
focused disease-management strategy \citep{2021-Phage_Ramos-Vivas}.
Hp3 differs from this benchmark by representing a deeply divergent
\emph{Aeromonas}-phage genus, but its approximately 45-minute latent
period and narrow host range (2 of 19 tested strains) should be viewed
as baseline biological properties rather than evidence of superior
biocontrol breadth \citep{2022-Characterization_Bujak,
2024-Novel_MuliyaSankappa}.

Although the full four-group survival design was statistically analysed,
several limitations temper the \emph{in vivo} interpretation. No
a priori power calculation was performed, and the randomization sequence,
allocation concealment, survival-monitoring blinding, and histopathology
assessment blinding were not formally documented. The animal challenge
was a single, unreplicated cohort, leaving possible cohort or tank effects
unresolved. Histopathology was descriptive and based on representative
images rather than blinded semi-quantitative lesion scoring, and tissue
bacterial loads were not measured. Finally, treatment was evaluated only
as a single intraperitoneal phage dose one hour after bacterial challenge,
while \emph{in vitro} suppression was compared at a single 48-h endpoint.

The characterization of Hp3 therefore yields a genetically distinctive
candidate rather than a finished field-ready product. Its complete
genome, divergent genus-level placement, lytic activity, and in vivo
rescue support further development, while resistance, host-range breadth,
delivery route, and mechanism-of-clearance questions define the next
experimental priorities.

\section{Conclusion}\label{conclusion}

In this study, we isolated and characterized Hp3, a lytic phage targeting
the aquaculture pathogen \emph{A. hydrophila}. Genomic analyses support
its placement in a deeply divergent \emph{Aeromonas}-phage genus within
the \emph{Caudoviricetes}, with no known antimicrobial-resistance genes
or high-confidence virulence determinants. Hp3 administration increased
the survival rate from 15\% to 85\% in a largemouth bass infection
model, demonstrating protective efficacy under the tested experimental
conditions. These findings position Hp3 as a promising but early-stage
antimicrobial candidate; future work should validate scalable delivery
routes, cocktail formulation, resistance management, and mechanisms of
in vivo bacterial clearance.

\section{Declarations}\label{declarations}

\subsection{Ethics statement}\label{ethics-statement}

All animal experiments were conducted in strict accordance with the
animal welfare and ethics guidelines established by the Institutional
Animal Care and Use Committee (IACUC) of Northwest A\&F University,
China. The specific experimental protocol for this study was reviewed
and approved by the IACUC (Approval Date: March 13, 2025). As per the
institutional practice, this project is authorized under the
Institutional License Number NWAFU-314020038.

\subsection{Competing interests}\label{competing-interests}

The authors declare that they have no competing interests.

\subsection{Funding}\label{funding}

This work was supported by the National Key Research and Development
Program of China (No.~2023YFD2400700) awarded to GW, the Deutsche
Forschungsgemeinschaft (DFG, German Research Foundation - Emmy Noether
program, Project No.~273124240; SPP2330, Project No.~464797012), and the
European Research Council Starting grant (ERC StG 803077) awarded to LD.

\subsection{Data availability}\label{data-availability}

All raw sequencing data and assembled genomes generated in this project
can be found in the NCBI BioProject database under the accession number
PRJNA1348731; the complete genome of Hp3 is available under GenBank
accession PX754746.1. The supplementary dataset supporting the findings of this
study, including the raw numerical data for all figures, is openly
available in Figshare under the Creative Commons Attribution 4.0 (CC-BY
4.0) license at: \url{https://doi.org/10.6084/m9.figshare.30500525}. The
source code used for data analysis is maintained at
\url{https://github.com/rujinlong/HydrophilaPhage} and permanently
archived in Zenodo, licensed under the GNU General Public License v3.0
(GPL-3.0), at: \url{https://doi.org/10.5281/zenodo.17454899}.

\subsection{Authors' contributions}\label{authors-contributions}

All authors read and approved the final manuscript.

Methodology: CL, YZ and JR; Investigation: CL, YZ, JX and HZ; Formal
analysis: CL, YZ, JX and JR; Visualization: CL, YZ and JR;
Conceptualization: FL, GW, and LD; Writing -- original draft: CL and JR;
Writing -- review \& editing: FL, GW and LD; Supervision: JX, GW and LD;
Funding acquisition: GW and LD.

\subsection{Acknowledgements}\label{acknowledgements}

The authors sincerely thank the members of the Deng laboratory for their
constructive discussions and valuable suggestions. CL and YZ sincerely
thank the Northwest A\&F University for awarding the Doctoral Students'
Internationalization Training Project.


\renewcommand\refname{Reference}


\bibliography{bibliography}

\end{document}
